\begin{proof}
\begin{enumerate}[label=\arabic*.]
\item We first fix the dual convention.  The preceding construction
has cohomology monodromies $A,B$; Theorem \ref{maps:src:F.2.5} uses homology
monodromies $A^t,B^t$.  All three special values are real.
Choose a positive real basepoint and paths along the positive
real circle to infinity and to zero.  Complex conjugation fixes
these paths and reverses the respective peripheral meridians.
Pull back the preceding period family by conjugation of the
base and conjugate its complex fibres.  Its Hodge subbundle is
holomorphic, with local expression
$\overline{\mathcal F(\bar\lambda)}$, and its cohomology
monodromies are $A^{-1},B^{-1}$.  Its homology monodromies are
therefore exactly $A^t,B^t$.
Put $M_1=A^t$, $M_2=B^t$, $M_o=M_2M_1$, $M_0=(M_1M_2)^{-1}$ and
$J=\operatorname{diag}(I_3,-I_3)$.  The based peripheral word is
$\kappa(g_0)=g_1^{-1}g_0^{-1}g_1$, so
$M_0=M_1M_o^{-1}M_1^{-1}$.  For the specified target marking $S$,
direct integral multiplication gives
\[
 S^{-1}M_0S=\begin{pmatrix}I_3&I_3\\0&I_3\end{pmatrix},\qquad
 S_o=M_1^{-1}SJ,\qquad
 S_o^{-1}M_oS_o=\begin{pmatrix}I_3&I_3\\0&I_3\end{pmatrix}.
\]
Any two markings of this positive standard form differ by
$G=\left(\begin{smallmatrix}P&Q\\0&P\end{smallmatrix}\right)$,
where $P\in\operatorname{GL}_3(\mathbb Z)$ and $Q$ is integral:
this follows by multiplying $G$ with the displayed unipotent
matrix.  The actual period change is
\[
 (I_3,Z)G=(P,Q+ZP),\qquad
 P^{-1}(P,Q+ZP)
  =(I_3,sI_3+P^{-1}Q+P^{-1}CP).
\]
Thus the bounded holomorphic logarithmic form holds in $S_o$.
In the real-normalized cusp coordinate,
$q_o=\bar q$ and $s_o=-\bar s$.
Base and fibre conjugation therefore give
$(I_3,-sI_3+C^\kappa(q))$, with
$C^\kappa(q)=\overline{C(\bar q)}$.
The period-cycle change $J$ gives the target marking
$M_1S_oJ=S$ and the positive periods
$(I_3,sI_3-C^\kappa(q))$.


\item For a holomorphic Hodge subbundle $\mathcal F$ with
$\mathcal F\cap\bar{\mathcal F}=0$, restriction of integral
homology functionals embeds the homology lattice discretely in
$\mathcal F^*$.  Real transversality makes its real span the
whole real vector space.  The quotient
$\mathcal F^*/\mathcal V_{\mathbb Z}^*$ is consequently a
holomorphic family of compact complex three-tori over $U$,
with its zero section.  This also proves local properness,
by a smoothly varying fundamental parallelepiped.

\item On the root disc at the finite point of order $m_i$, the family
extends smoothly and its deck map has linear part $M_i$.
The invariant vector $v_i$ of Theorem \ref{maps:src:F.2.5} gives a holomorphic
torsion section $v_i/m_i$.  Use the diagonal affine deck action
\[
 (z,s)\longmapsto(M_i z+v_i/m_i,\zeta_{m_i}s).
 \tag{F.24}\label{maps:eq:F.24}
\]
Here $M_i z$ denotes the induced complex-linear map
$\mathcal F_s^*\to\mathcal F_{\zeta_{m_i}s}^*$ from the marked
lattice action, not multiplication of a complex three-vector
by a six-by-six matrix.  Likewise $v_i/m_i$ means its evaluation
as a point of the target torus.
It has order $m_i$, since the $m_i$th translation is integral.
Every nonidentity power is fixed-point-free on the central
torus: the invariant integral covector $\ell$ takes $v_i$ to
$1$ or $-1$, so a fixed point would require the noninteger
$\pm j/m_i$ to vanish in $\mathbb R/\mathbb Z$.
Away from the central fibre the base rotation is already free.
The quotient is therefore smooth and proper over the coarse
disc, with precisely the required multiple fibre.

\item Its marked punctured family is holomorphically identified with
the original torus family by logarithmic translation.  On the
root cover use $r(s)=(\log s/(2\pi i))v_i$.  Then
$r(\zeta_{m_i}s)-M_i r(s)=v_i/m_i$, modulo the integral lattice;
the change of branch of $\log s$ is itself an integral vector.
Thus translation by $r$ conjugates the linear deck action to
\eqref{maps:eq:F.24}.  This fixes the affine clutching, including its sign.

\item At the cusp use the integral marking of \eqref{maps:eq:F.13}, the fan \eqref{maps:eq:E.4}
and the action \eqref{maps:eq:E.5} with $C_3=C$.
Its maximal charts are $\mathbb C^4$ with map
$t=z_0z_1z_2z_3$, and its transition maps are the toric changes
of coordinates multiplied by the holomorphic units
$\exp(2\pi i C(t)\lambda)$.
The actual multiplicative comparison is
\[
 q=\bar q_o,\qquad x=(\bar x_o)^{-1},\qquad
 \nu=-\nu_o,
 \qquad C_{\rm new}(q)=-C^\kappa(q).
\]
It carries the old period action to
$x\mapsto\exp(2\pi iC_{\rm new}(q)\nu)q^\nu x$.
On cocharacters it is
$\Phi=\operatorname{diag}(-I_3,1)$ together with conjugation.
The Kuhn simplex at $v$ with ordering $\sigma$ is carried to
the simplex at $-v-(1,1,1)$ with reversed ordering.
If $N_o$ and $N$ are their height-one generator matrices and
$R$ reverses four columns, then
$N=\Phi N_oR$ and $N^{-1}\Phi=RN_o^{-1}$.
Consequently the full affine-chart map is
$z_j=\overline{z_{o,3-j}}$.  These character identities also
give equality on every monomial overlap; the product of the
four coordinates is the conjugate base coordinate.
The already displayed period-action equality makes the map
descend to the analytic period quotient.  Hence the fan and
marking in this comparison are precisely the prescribed ones.

The proof of Theorem \ref{maps:src:E.1.2} used only boundedness of $C$,
integrality of the standard cusp, and unimodularity of that
fan.  It therefore proves smoothness, proper discontinuity
and properness for this model without a symmetry condition
on $C$.  Glue by the displayed integral marking and use no
additional cusp translation.

\item The three punctured overlap maps are holomorphic isomorphisms
with the same marked torus family, so their cocycle condition
is literal.  Proper local models and a compact complementary
base give a Hausdorff compact total space.  Every chart is
smooth.  The affine vectors, cusp fan and clutching are exactly
those in Theorems \ref{maps:src:F.2.5}--\ref{maps:src:F.2.6}.  Varying their bounded corrections
and varying the fibre complex structures along a smooth path
does not change these marked smooth fillings: on the finite
root models the underlying real torus family is the fixed
local-system bundle; on the cusp use the stratum-preserving
parameter vector fields and uniform properness estimate of
Theorem \ref{maps:src:E.1.3}.  Thus the resulting underlying smooth manifold
is the one whose integral topology was computed there.

\end{enumerate}
\end{proof}
